% ============================================================================
%  Chapter 7. Evaluation: Four Independent Measurements of Failure
%
%  \input file. Requires booktabs, amsmath, natbib in the parent preamble.
%  Bibliography: accord/thesis/refs.bib
%
%  Traceability: the audit scripts over the 299 pharmacist-annotated pairs,
%  transfer_results_correct.json, the MedNLI fine-tuning logs, and
%  accord/CORRECTIONS.md (C2, C4, C5).
% ============================================================================

\chapter{Evaluation: Four Independent Measurements of Failure}
\label{ch:eval}

\section{Evaluation Design}
\label{sec:eval-design}

A negative result is only worth reporting if it is difficult to escape. A single
unfavourable metric invites the response that the metric was the wrong one, the
threshold was badly chosen, the class balance was hostile, or the sample was too
small. This chapter is constructed to make those responses unavailable, by
measuring the same system four ways that fail for different reasons if they fail
at all.

\begin{enumerate}
\item \textbf{Ranking.} Does the model's contradiction probability order the
      pharmacist's labels better than chance? This is threshold-free and
      insensitive to class balance. Section~\ref{sec:m1}.
\item \textbf{Operating points.} Is there \emph{any} threshold at which the
      model is usable? This removes the objection that $0.5$ was the wrong
      cut. Section~\ref{sec:m2}.
\item \textbf{Severity.} Does the model's harm score correlate with the
      pharmacist's $1$ to $5$ danger rating? This asks a different question of a
      different output head. Section~\ref{sec:m3}.
\item \textbf{Comparison.} How does the model rank against a general-purpose
      language model, zero-shot and few-shot, and against the weak labeller it
      was distilled from, on the same $299$ pairs? This removes the objection
      that the task is simply impossible. Section~\ref{sec:m4}.
\end{enumerate}

All four are computed on the $299$ pharmacist-annotated pairs, which are absent
from the training set by construction (Chapter~\ref{ch:weak},
Section~\ref{sec:weak-setup}). The single configuration that trained on those
pairs is not reported anywhere in this thesis.

Section~\ref{sec:eval-scoping} then addresses the obvious objection: given the
exemplar defect of Chapter~\ref{ch:candidates}, can any of this be believed? The
answer is a qualified yes, and the qualification is stated precisely rather than
waved at.

\section{A Component Sanity Check: BioMedBERT on MedNLI}
\label{sec:mednli}

\subsection{Purpose}

Before reporting that a system fails, it is worth establishing that the pieces
were assembled competently. MedNLI \citep{romanov2018mednli} is the standard
clinical natural language inference benchmark: premise and hypothesis pairs derived
from MIMIC-III, labelled entailment, contradiction, or neutral. It is the
closest well-established task to the one this thesis attempts, and it has an
established performance range.

\subsection{Result}

\begin{table}[t]
\centering
\small
\begin{tabular}{lrrr}
\toprule
\textbf{Configuration} & \textbf{Trainable} & \textbf{Test acc.} & \textbf{Macro $F_1$} \\
\midrule
BioMedBERT, full fine-tune, 10 epochs        & full backbone & 83.90\% & 83.91\% \\
BioMedBERT $+$ LoRA ($r{=}16$, $\alpha{=}32$), 5 epochs & 592,131 & 72.93\% & 72.74\% \\
\midrule
Published range for comparable systems       & n/a & \multicolumn{2}{c}{83 to 87\%} \\
\bottomrule
\end{tabular}
\caption{MedNLI results. The full fine-tune sits inside the published range for
comparable encoder-based systems. This establishes that the backbone, the
tokenisation, the training loop, and the evaluation harness are sound. The LoRA
configuration used for ACCORD loses roughly eleven accuracy points, which is a
material cost that was not accounted for in the project's design.}
\label{tab:mednli}
\end{table}

\subsection{What this row is, and what it is not}

\textbf{Table~\ref{tab:mednli} is a BioMedBERT baseline. It is not ACCORD.} Zero
ACCORD components are present in either model: no cross-note encoder, no
authority conditioning, no prior conditioning, no harm estimation. It is a
sentence-pair classifier on a public benchmark.

This distinction is emphasised because the alternative framing, reporting
$83.90\%$ under the project's own system name, would be the single most
misleading move available in this thesis. A reader encountering
``ACCORD: $83.9\%$ on MedNLI'' followed by the results in
Section~\ref{sec:eval-point} would reasonably conclude that a strong model
failed on a hard task. What actually happened is that a strong \emph{backbone}
performs as expected on a benchmark it was designed for, and a differently
constructed system built on that backbone fails on a different task for reasons
that have nothing to do with the backbone's capability.

\subsection{One substantive finding}

The eleven-point gap between full fine-tuning and LoRA is worth recording. LoRA
is often described as approximately loss-free relative to full fine-tuning
\citep{hu2022lora}, and on many tasks it is. Here, at $r = 16$ with $592{,}131$
trainable parameters and a shorter schedule, it is not. ACCORD used the LoRA
configuration. Whether the gap reflects rank, schedule length, or the specific
demands of clinical inference is untested here. [TODO: run a rank sweep
$r \in \{8, 16, 64, 256\}$ at matched epochs on MedNLI and report the
accuracy and parameter frontier. The two rows in Table~\ref{tab:mednli} differ in
both rank and schedule and therefore do not isolate either.]

\section{Point Metrics on the Gold Set}
\label{sec:eval-point}

\begin{table}[t]
\centering
\small
\begin{tabular}{lr}
\toprule
\textbf{Quantity} & \textbf{Value} \\
\midrule
Annotated pairs                                   & 299 \\
\quad true contradictions                         & 12 \\
\quad not contradictions                          & 241 \\
\quad ambiguous                                   & 46 \\
\midrule
Pairs ACCORD predicts \emph{contradiction} at $\tau = 0.5$ & 278 \\
\midrule
Precision                                         & 0.047 \\
Recall                                            & 0.917 (11/12) \\
Accuracy                                          & 0.115 \\
Pooled base rate                                  & 0.040 (12/299) \\
\bottomrule
\end{tabular}
\caption{Point metrics at the default threshold. The model predicts
\emph{contradiction} for $93\%$ of the evaluation set. Accuracy of $0.115$ on a
task whose majority class is $80.6\%$ of the labelled pairs is the accuracy of a
near-constant positive predictor evaluated against a heavily negative
population. The precision of $0.047$ should be read against the base rate rather
than in isolation.}
\label{tab:point}
\end{table}

\paragraph{On denominators.} Precision and accuracy in Table~\ref{tab:point} are
computed on the $253$ non-ambiguous pairs, while the predicted-positive count of
$278$ is over all $299$.\footnote{The reconstruction is as follows. The
threshold sweep in Section~\ref{sec:m2} reports precision $0.0474$ at a
positive rate of $1.0$, and $12/253 = 0.0474$, which fixes the denominator at
$253$. With $11$ true positives, a precision of $0.047$ requires $234$ predicted
positives among those $253$, and an accuracy of $0.115$ then follows as
$29/253$. The $278$ figure over all $299$ implies that $44$ of the $46$
ambiguous pairs are also predicted positive. [TODO: confirm these denominators
directly in the audit script and report them explicitly in the results file. The
present mixture of $n{=}253$ and $n{=}299$ denominators in a single table is a
reporting defect, corrected here by annotation rather than by recomputation.]}
On the $253$-pair basis the comparable base rate is $12/253 = 4.74\%$, and
ACCORD's precision of $4.70\%$ is very slightly \emph{below} it.

\paragraph{The reading.} A detector whose precision equals the base rate to
within four hundredths of a percentage point is not a detector. It is a
relabelling of the candidate pool. The recall of $0.917$ is real, but it is the
recall that any near-constant positive predictor achieves: predicting
\emph{contradiction} for everything yields recall $1.0$ at precision equal to
the base rate, and ACCORD is a small perturbation of that predictor.

This is the point at which most evaluations stop, and it is also the point at
which the result is most easily dismissed. The next four sections make it
undismissable.

\section{Measurement 1: The Score Carries No Ranking Information}
\label{sec:m1}

\subsection{Why ranking is the decisive test}

Precision and recall depend on a threshold, and thresholds can be argued about.
Area under the ROC curve does not. It is the probability that a randomly drawn
positive receives a higher score than a randomly drawn negative. If the model
has learned \emph{anything} about the task, even something too weak to support a
usable threshold, AUROC exceeds $0.5$. If AUROC is $0.5$, the score is
statistically independent of the label.

\begin{table}[t]
\centering
\small
\begin{tabular}{llrr}
\toprule
\textbf{Score} & \textbf{Label treatment} & \textbf{$n$} & \textbf{AUROC} \\
\midrule
\texttt{accord\_contra\_prob} & all pairs, ambiguous excluded from positives & 299 & 0.460 \\
\texttt{accord\_contra\_prob} & ambiguous pairs dropped                       & 253 & 0.460 \\
\texttt{accord\_contra\_prob} & ambiguous counted as positive                 & 299 & 0.490 \\
\midrule
\texttt{accord\_harm\_score}  & all pairs                                     & 299 & 0.504 \\
\bottomrule
\end{tabular}
\caption{Ranking performance under three treatments of the ambiguous class. The
result is stable at approximately $0.46$ to $0.49$ regardless of how ambiguity is
handled, which forecloses the objection that the ambiguous stratum is driving
the finding. The harm score, which Chapter~\ref{ch:arch} established was never
trained, sits at $0.504$: the value expected of a randomly initialised head.}
\label{tab:auroc}
\end{table}

\subsection{The result}

Table~\ref{tab:auroc} reports AUROC of $0.460$ for the contradiction
probability. This is below chance, and the supporting statistics establish that
it is indistinguishable from chance rather than meaningfully inverted:

\begin{itemize}
\item \textbf{Average precision} between $0.039$ and $0.046$ across the three
      label treatments, against a base rate of $0.040$. Average precision is the
      appropriate summary under this degree of class imbalance
      \citep{saito2015prcurve}, and it sits exactly at the no-skill line.
\item \textbf{Mann-Whitney $U$} comparing the score distributions of positives
      and negatives: $p = 0.64$. The two distributions are not separated.
\item \textbf{Bootstrap $95\%$ confidence interval} on AUROC, $B = 10{,}000$
      resamples \citep{efron1994bootstrap}: $[0.287, 0.633]$. The interval
      contains $0.5$ comfortably. It is also very wide, which is the honest
      consequence of $12$ positives.
\end{itemize}

\subsection{The top-$k$ check}

AUROC is a global summary and can obscure behaviour at the extremes, which is
where a triage system would actually operate. \textbf{The $15$ highest-scored
pairs in the evaluation set contain zero true contradictions.} Under the null
hypothesis of random ordering, the expected count is $15 \times 12/299 = 0.6$,
so observing zero is unremarkable in itself. That is precisely the point: the
model's most confident predictions are indistinguishable from a random draw.

For a deployed triage system this is the operationally decisive number. The
entire value proposition of such a system is that a clinician reviews the top of
the ranked list. Here the top of the list is noise.

\subsection{The harm score}

The harm score's AUROC of $0.504$ at $p = 0.96$ is reported for completeness and
should be interpreted as Chapter~\ref{ch:arch}, Section~\ref{sec:c3} instructs:
this head received no loss term at any point in training. Its output range
across the $299$ pairs is $0.398$ to $0.610$, the range of an untrained sigmoid.
A null result here is a measurement of an untrained module, not evidence about
whether clinical harm is predictable.

\section{Measurement 2: No Operating Point Exists}
\label{sec:m2}

\subsection{The objection this forecloses}

The natural response to Section~\ref{sec:eval-point} is that $\tau = 0.5$ is the
wrong threshold for a $4\%$-prevalence task and that a lower threshold, or a
prior correction, would produce a usable system. This section shows that no
threshold produces a usable system, because there is nothing in the score
ordering for a threshold to exploit.

\subsection{The sweep}

\begin{table}[t]
\centering
\small
\begin{tabular}{lrrr}
\toprule
\textbf{Threshold $\tau$} & \textbf{Precision} & \textbf{Recall} & \textbf{Positive rate} \\
\midrule
0.05 & 0.0474 & 1.000 & 1.000 \\
0.10 & 0.0474 & 1.000 & 1.000 \\
0.15 & 0.0474 & 1.000 & 1.000 \\
\bottomrule
\end{tabular}
\caption{Threshold sweep. At every threshold in the low range the model predicts
positive for \emph{every} pair. The precision of $0.0474$ is exactly $12/253$,
the base rate: at a positive rate of $1.0$ precision is arithmetically forced to
the base rate. The three rows are identical because no score in the evaluation
set falls below $0.15$.}
\label{tab:sweep}
\end{table}

\begin{table}[t]
\centering
\small
\begin{tabular}{lr}
\toprule
\textbf{Target recall} & \textbf{Achieved precision} \\
\midrule
0.500 & 5.4\% \\
0.750 & 4.1\% \\
0.917 & 4.1\% \\
\bottomrule
\end{tabular}
\caption{Precision at fixed recall. Even at $50\%$ recall, precision reaches
only $5.4\%$: roughly eighteen false alarms for every true contradiction. The
precision-recall curve is essentially flat at the base rate, which is the
signature of a score that does not order the labels. [TODO: the $4.1\%$ figure
at recall $0.917$ and the $4.70\%$ reported at $\tau = 0.5$ in
Table~\ref{tab:point} are computed on different denominators ($n = 299$ against
$n = 253$); reconcile and report both on a single basis.]}
\label{tab:pr}
\end{table}

\subsection{The reading}

Tables~\ref{tab:sweep} and~\ref{tab:pr} together establish that the model's
scores occupy a narrow band above $0.15$ and that within that band the ordering
is uninformative. No threshold, no prior correction, and no post-hoc calibration
recovers a usable operating point, because calibration reshapes a score
distribution and cannot create an ordering that is not there
\citep{guo2017calibration}.

This is the answer to a question Chapter~\ref{ch:weak} deferred. The prevalence
mismatch in the weak labels could in principle have been corrected after the
fact by shifting the decision threshold. It could not be corrected here, because
the damage was not a shifted threshold on an informative score. It was the loss
of the score's informativeness altogether.

\section{Measurement 3: Severity Ranking}
\label{sec:m3}

The pharmacist assigned a $1$ to $5$ danger rating to each pair. Two Spearman rank
correlations were computed against it:

\begin{table}[h]
\centering
\small
\begin{tabular}{lrrr}
\toprule
\textbf{Model score} & \textbf{Spearman $\rho$} & \textbf{$p$} & \textbf{$n$} \\
\midrule
\texttt{accord\_harm\_score}  & $-0.0001$ & 0.999 & 299 \\
\texttt{accord\_contra\_prob} & $-0.029$  & 0.62  & 299 \\
\bottomrule
\end{tabular}
\caption{Severity ranking. Both correlations are null. The harm score's
$\rho = -0.0001$ at $p = 0.999$ is about as close to an exact null as a finite
sample produces.}
\label{tab:severity}
\end{table}

\paragraph{How to read this, and how not to.} The harm-score row is \emph{not}
evidence that clinical severity is unpredictable. Chapter~\ref{ch:arch},
Section~\ref{sec:c3} established that the harm head receives no loss term, its
guideline bank is twenty random vectors, and its outputs span the range of an
untrained sigmoid. A correlation of $-0.0001$ between a randomly initialised
head's outputs and a clinical rating is the expected result and confirms the
architectural finding.

\textbf{The severity-ranking research question was therefore not tested by this
work.} It remains open. What this thesis establishes about it is narrower and
still useful: a severity head can be specified, implemented, run end to end,
included in a results table, and reported as a contribution without ever
appearing in the loss, and the only measurement that catches this is the
empirical range of its outputs.

The contradiction-probability row is a genuine, if unsurprising, null: a score
that does not separate contradictions from non-contradictions
(Section~\ref{sec:m1}) has no reason to rank them by danger.

\section{Measurement 4: System Comparison}
\label{sec:m4}

\subsection{The objection this forecloses}

A failure that is universal across methods is a statement about the task, not
about the system. If every approach performs at the base rate, the correct
conclusion is that cross-note contradiction detection at $4\%$ prevalence from a
two-sentence view is not solvable, and the specific defects catalogued in
Chapters~\ref{ch:candidates} through~\ref{ch:arch} are beside the point.

\begin{table}[t]
\centering
\small
\begin{tabular}{lrrr}
\toprule
\textbf{System} & \textbf{Precision} & \textbf{Recall} & \textbf{MCC} \\
\midrule
GPT-4o, zero-shot        & \textbf{8.45\%} & 50.0\% & $\mathbf{+0.109}$ \\
GPT-4o, few-shot         & 4.96\%          & 58.3\% & $+0.012$ \\
ACCORD, $\tau = 0.5$     & 4.70\%          & 91.7\% & $-0.007$ \\
Llama 3 8B (weak labeller) & 4.55\%        & 91.7\% & $-0.044$ \\
\bottomrule
\end{tabular}
\caption{System comparison on the $299$ pharmacist-annotated pairs. Matthews
correlation coefficient \citep{matthews1975mcc} is reported because it is the
summary statistic least distorted by extreme class imbalance; a value of $0$ is
chance. GPT-4o zero-shot is the only system with a positive MCC. ACCORD is
indistinguishable from chance. The weak labeller ACCORD was distilled from is
\emph{worse} than chance.}
\label{tab:systems}
\end{table}

\subsection{The four findings in Table~\ref{tab:systems}}

\paragraph{The task is not impossible.} GPT-4o zero-shot
\citep{openai2023gpt4} achieves $8.45\%$ precision at $50\%$ recall with a
positive MCC. This is not a usable clinical system, and it should not be
presented as one: eleven false alarms per true contradiction, and half the true
contradictions missed. But it is above chance, on the same evidence, on the same
$299$ pairs. Whatever is wrong with ACCORD is not that the evidence contains no
signal.

\paragraph{Few-shot prompting made GPT-4o worse.} Precision falls from $8.45\%$
to $4.96\%$ and MCC from $+0.109$ to $+0.012$, while recall rises from $50.0\%$
to $58.3\%$. The exemplars pulled the model towards the positive class. This is
consistent with the class-prior sensitivity documented for in-context learning
\citep{zhao2021calibrate}, and it is the same failure direction as the weak
labeller's, arrived at independently. [TODO: report the exemplar class balance
used in the few-shot prompt. If it over-represents positives relative to the
$4\%$ base rate, the mechanism is identified rather than hypothesised.]

\paragraph{ACCORD did not improve on its teacher, and the teacher was worse than
chance.} ACCORD ($-0.007$) and Llama ($-0.044$) sit at and below zero
respectively, with identical recall of $91.7\%$ and near-identical precision.
This is Chapter~\ref{ch:weak}'s finding restated on the gold set: the student
reproduced the teacher.

\paragraph{The purpose-built system is the second-worst of four.} A model
carrying roughly $7.5$M task-specific parameters over a domain-pretrained
biomedical encoder, trained on $1{,}679$ in-domain examples, performs worse on
the target task than a general-purpose model given a single prompt and no
training data at all.

\section{Error Taxonomy}
\label{sec:errors}

\subsection{What the twelve true contradictions are}

\begin{table}[t]
\centering
\small
\begin{tabular}{lrr}
\toprule
\textbf{Contradiction type} & \textbf{Count} & \textbf{Share} \\
\midrule
Omission (drug present in one note, absent from the other) & 7 & 58.3\% \\
Dosage discrepancy                                          & 3 & 25.0\% \\
Drug substitution                                           & 1 & 8.3\% \\
Class conflict                                              & 1 & 8.3\% \\
\midrule
\textbf{Total}                                              & \textbf{12} & \\
\bottomrule
\end{tabular}
\caption{Types of the twelve pharmacist-confirmed contradictions. Omission
dominates. This matches the medication reconciliation literature, where omission
of a regularly used medication is the single most common category of unintended
discrepancy \citep{cornish2005discrepancies}, and it means the hardest evidential
case, reasoning about something that is \emph{not} in a document, is also the
most clinically important one.}
\label{tab:errtypes}
\end{table}

\begin{table}[t]
\centering
\small
\begin{tabular}{lrrrrr}
\toprule
\textbf{Danger rating} & \textbf{1} & \textbf{2} & \textbf{3} & \textbf{4} & \textbf{5} \\
\midrule
Count & 2 & 4 & 2 & 1 & 3 \\
\bottomrule
\end{tabular}
\caption{Pharmacist danger ratings of the twelve true contradictions. Four of
twelve are rated $4$ or $5$. These are the cases a triage system exists to
surface, and Section~\ref{sec:m3} established that no model score orders them.}
\label{tab:danger}
\end{table}

\subsection{Who catches what}

\begin{table}[h]
\centering
\small
\begin{tabular}{lr}
\toprule
\textbf{System} & \textbf{True contradictions recovered} \\
\midrule
Llama 3 8B          & 11 / 12 \\
ACCORD              & 11 / 12 \\
GPT-4o, zero-shot   & 6 / 12 \\
\bottomrule
\end{tabular}
\caption{Recall on the twelve true contradictions. The two near-constant
positive predictors recover eleven; the selective predictor recovers six. This
table should not be read as evidence that the first two are better detectors,
and Section~\ref{sec:errors-reading} explains why.}
\label{tab:recall}
\end{table}

\subsection{How to read Table~\ref{tab:recall}}
\label{sec:errors-reading}

Recall in isolation is the most misleading metric available in this setting. A
system that predicts \emph{contradiction} for everything achieves recall
$1.0$. ACCORD and Llama achieve $11/12$ because they predict positive for the
overwhelming majority of pairs: ACCORD for $278$ of $299$, and Llama at a
comparable rate, consistent with the $84.9\%$ positive rate
Chapter~\ref{ch:weak} measured over the full candidate pool. GPT-4o achieves
$6/12$
while predicting positive far less often, and consequently is the only system
with a positive MCC.

The one substantive observation available from Table~\ref{tab:recall} concerns
ACCORD's single miss. Chapter~\ref{ch:candidates}, Section~\ref{sec:gold},
established arithmetically that ACCORD recovers one of the two true Rule 4
contradictions, and Rule 4 is the only rule on which the model exercises
meaningful selectivity ($74.5\%$ positive rather than $\approx 97\%$). The one
place the model declines to predict positive is the one place it declines
incorrectly. There is no configuration of this model in which its selectivity
helps.

\subsection{The clinically important observation}

Omission is $7$ of $12$ true contradictions and it is the category the
architecture is least equipped for. Detecting an omission requires reasoning
about absence, and the model's input for a Rule 4 candidate is a nursing
sentence paired with an empty string. There is no discharge evidence to encode,
so an asymmetric cross-note encoder has nothing to attend to on one side, and
the pair representation degenerates to a single-sentence representation. A
system designed around cross-note attention is structurally poorly matched to
the majority error type in its own target population. This is a design finding
independent of every implementation defect in Chapter~\ref{ch:arch}, and it
would apply equally to a correct implementation of the same specification.

\section{What the Exemplar Defect Does and Does Not Contaminate}
\label{sec:eval-scoping}

The most serious objection to this chapter is that Chapter~\ref{ch:candidates}
established that $41.4\%$ of action-conflict candidates were annotated against
evidence that did not carry their own conflict. If the labels are contaminated,
are the measurements meaningless? The answer requires separating three distinct
claims.

\subsection{What is genuinely bounded}

The claim ``the candidate pool has a precision of $4.0\%$'' is an
\emph{upper bound on the validity of the annotation exercise}, not an estimate
of pipeline quality. The pharmacist rendered correct judgements about the
evidence presented; for a large minority of action-conflict candidates that
evidence was not the evidence the rule had found. We do not know how the same
annotator would have ruled on corrected exemplars, and both a higher and a
lower true precision are consistent with the data collected. Every use of the
$4.0\%$ figure in this thesis carries this qualification.

\subsection{What is not contaminated}

The four measurements in this chapter are statements about whether the model's
outputs carry information about the pharmacist's labels. Both sides of each
comparison refer to \emph{the same object}: the model scored exactly the
exemplar sentence pairs the pharmacist annotated. The question ``does the model's
score order the pharmacist's label'' is therefore well-posed regardless of
whether that label is a good proxy for clinical truth. If the model had learned
any function of the presented sentence pair, and the pharmacist's judgement is
any function of that same pair, the two would correlate. They do not.

The same argument covers the threshold sweep, the severity correlations, the
weak-label marginal of Chapter~\ref{ch:weak}, and the transfer collapse of
Chapter~\ref{ch:arch}. None of these compares a model output against a claim
about clinical reality; all compare model outputs against each other or against
annotations of the identical input.

\subsection{The residual risk, stated honestly}

There is a real residual risk and it should not be minimised. Label noise
attenuates a measured AUROC towards $0.5$. In principle, a model with a true
AUROC of, say, $0.65$ could be measured at $0.46$ under sufficient label noise.
Four considerations bound how much weight this deserves.

\begin{enumerate}
\item \textbf{Most of the evaluation set is unaffected.} The defect applies only
      to action-conflict candidates, Rules 1 and 2, which are $165$ of the $300$
      sampled pairs. Applying the pooled $41.4\%$ rate gives approximately $68$
      affected pairs, under a quarter of the evaluation set. Rules 3 and 4,
      $135$ pairs, are free of the defect by construction, because Rule 3
      selects the maximum-difference dose pair and Rule 4 has no discharge
      exemplar to mis-select.
\item \textbf{The decisive experiment is available and was not run.}
      [TODO: report AUROC, average precision, and the bootstrap interval
      restricted to the Rule 3 and Rule 4 strata, $n = 135$, which are free of
      the exemplar defect. If the restricted AUROC is also near $0.5$, the
      objection is closed. This is the single highest-value remaining
      measurement in this thesis and should be run before submission.]
\item \textbf{A near-constant predictor has no signal to attenuate.} The model
      predicts positive for $278$ of $299$ pairs and all its scores lie above
      $0.15$ in a narrow band. Label noise degrades an ordering; it does not
      explain the near-absence of an ordering, and Chapter~\ref{ch:weak}
      supplies an independent, sufficient explanation for that absence in the
      $21\times$ prevalence mismatch of the supervision.
\item \textbf{The top-$k$ result is robust to moderate label noise.} A model
      with genuine signal concentrates true positives at the top of its ranking
      even under substantial label corruption. Zero true contradictions in the
      top fifteen is not what an attenuated signal looks like.
\end{enumerate}

Point 2 is a genuine hole in the present work and is marked as such rather than
argued around.

\section{Synthesis}
\label{sec:synthesis}

\begin{table}[t]
\centering
\small
\begin{tabular}{p{0.3\linewidth}p{0.2\linewidth}p{0.42\linewidth}}
\toprule
\textbf{Stage} & \textbf{Status} & \textbf{Decisive measurement} \\
\midrule
Cohort construction & Works & $9{,}783$ admissions, $19{,}566$ notes, documented
retreat from three roles to two \\
\addlinespace
Medication extraction & Unvalidated & $8.9{:}1$ discharge/nursing asymmetry;
never annotated \\
\addlinespace
Candidate generation & Enriches, but evidence selection was wrong & Old exemplar
carried the conflict $58.6\%$ of the time; $739/1{,}307$ changed \\
\addlinespace
Weak supervision & Fails & $84.9\%$ positive against a $4.0\%$ base rate;
$21\times$ \\
\addlinespace
Architecture & Three claims falsified & Authority parameter moved $1.8\%$/$0.6\%$
from its initialisation; rule $F_1$ $0.96 \to 0.065$ under transfer; harm head
absent from the loss \\
\addlinespace
Detection & Fails & AUROC $0.460$, CI $[0.287, 0.633]$, $p = 0.64$; zero true
positives in the top fifteen \\
\addlinespace
Severity ranking & Not tested & $\rho = -0.0001$ from an untrained head \\
\bottomrule
\end{tabular}
\caption{Where the pipeline stands, stage by stage. Two stages work, one is
unvalidated, one is defective, and three fail. The failures are independent:
fixing the exemplar selector would not correct the supervision marginal, and
correcting the marginal would not connect the harm head to the loss.}
\label{tab:synthesis}
\end{table}

\subsection{The claim this thesis makes}

The standard pipeline for this problem, rule-based candidate mining, language
model weak supervision, and a purpose-built neural detector, fails at every
stage, and the failures are measurable and diagnosable rather than incidental.
Table~\ref{tab:synthesis} is that claim in full.

Three properties of the failure make it worth reporting rather than repairing
quietly.

\paragraph{It is not one failure.} The exemplar defect, the prevalence mismatch,
the initialisation artifact, the label leakage, and the untrained head are
causally independent. Each has its own mechanism and its own detector. A project
that fixed any one of them would still have produced a system with AUROC near
$0.5$.

\paragraph{Every failure was detectable before the results table.} Not one of
them required a held-out set, an expensive annotation, or a novel diagnostic.
The full detection kit is: a base-rate estimate before weak labelling, a printed
label marginal, a distinct-value count on the confidence field, a module
manifest with loss-graph membership, a parameter-displacement log, an
out-of-domain transfer run, an output-range check, a majority-class baseline,
and a group-disjointness assertion. Chapter~\ref{ch:arch},
Section~\ref{sec:diagnostics} states them as a protocol.

\paragraph{The negative result is measured four ways.} Ranking, thresholds,
severity, and comparison against three external systems. They fail for different
reasons if they fail, and they agree.

\subsection{What remains true after all of this}

The clinical problem is real. Omission of a regularly used medication is the
most common category of unintended discrepancy at care transitions
\citep{cornish2005discrepancies}, and adverse events following discharge are
common, frequently drug-related, and frequently preventable
\citep{forster2003adverse}. A board-certified pharmacist found $12$ genuine
contradictions in $299$ pairs, four of them rated $4$ or $5$ for danger. GPT-4o
zero-shot detects half of them above chance. Nothing in this thesis suggests the
problem is not worth solving. It establishes that this particular route to
solving it does not work, and it identifies, at each stage, the measurement that
would have said so earlier.

\section{Limitations}
\label{sec:eval-limitations}

\paragraph{Twelve positives.} Every confidence interval in this chapter is wide
because the evaluation set contains twelve positive examples. The bootstrap
interval on AUROC, $[0.287, 0.633]$, is consistent with a genuinely inverted
model and with a mildly informative one. What it is \emph{not} consistent with
is a useful one, and that is the claim being made. A larger annotated set is the
single most valuable extension of this work.

\paragraph{One annotator, no inter-rater reliability.}
Chapter~\ref{ch:candidates}, Section~\ref{sec:gold}. The previously reported
$\kappa = 1.000$ is retracted as circular. Inter-rater reliability is not
established for any label used in this chapter.

\paragraph{The exemplar defect.} Section~\ref{sec:eval-scoping}, including the
[TODO] that would close the objection.

\paragraph{Recall is measured only within the candidate pool.} Nothing outside
the $1{,}978$ candidates is detectable by any system evaluated here, including
GPT-4o, which was given the same candidates. The rules' recall against all true
contradictions in the $9{,}783$-admission cohort is unmeasured.

\paragraph{No efficiency analysis is reported.} Training cost, inference
latency, and the cost of the GPT-4o comparison were not systematically recorded.
[TODO: report wall-clock training time, inference latency per pair, and total
API cost for the GPT-4o rows. For a system whose value proposition is triage at
scale, cost per screened admission is a first-order metric and its absence is a
gap.]

\paragraph{One backbone, one adapter configuration.} Every ACCORD result uses
BioMedBERT with LoRA at $r = 16$. Section~\ref{sec:mednli} shows this
configuration costs eleven accuracy points against full fine-tuning on MedNLI.
Whether a full fine-tune of the same architecture would behave differently is
untested. It would not repair the supervision, but the claim that it changes
nothing is not established. [TODO: retrain ACCORD with a full backbone
fine-tune on the identical held-out split and report AUROC.]

\paragraph{Single institution.} MIMIC-III only. The transfer results in
Chapter~\ref{ch:arch} move across index diagnosis, not across institution.

\section{Ethics Statement}
\label{sec:eval-ethics}

\paragraph{Data governance.} MIMIC-III is credentialed data used under a data
use agreement. No note text appears in this thesis or in any derived artifact,
all note-derived intermediates are confined to a gitignored local directory, and
no re-identification was attempted. The one open governance question concerns
the GPT-4o comparison in Section~\ref{sec:m4}. [TODO: state the governance basis
under which candidate excerpts were transmitted to a hosted API, including
whether zero-retention terms were in force. If this cannot be established, the
GPT-4o rows should carry that caveat wherever they are reported.]

\paragraph{Deployment risk of the system as measured.} A system with $91.7\%$
recall and $4.7\%$ precision deployed as a clinical alert generates roughly
twenty false alarms per true contradiction. The established consequence of alert
burdens at that ratio is desensitisation, and desensitisation degrades response
to the alerts that matter. \textbf{This system should not be deployed in any
form.} The stronger statement is that its failure mode is worse than doing
nothing, because it would consume clinician attention while providing, by
Section~\ref{sec:m1}, no information.

\paragraph{Automation bias.} A component of this work presented a severity score
to a downstream consumer. Chapter~\ref{ch:arch} established that the score comes
from an untrained head. A clinical decision support system displaying a
plausible-looking numeric danger score derived from random initialisation is a
specific and serious hazard, because numeric scores in clinical interfaces are
acted upon. That this reached a results table is the most serious individual
failure documented in this thesis, and the diagnostic that catches it, printing
the empirical range of every emitted score, costs one line.

\paragraph{Publication ethics.} Four claims that appeared in an earlier
submission of this work are retracted here: the discovered authority hierarchy,
the rule-type $F_1$ of $0.9603$, the internal-consistency $\kappa$ of $1.000$,
and the intermediate-documentation prefilter finding. Two further reporting
defects are corrected: results from a configuration trained on the evaluation
set are withdrawn, and an ablation table reporting three seeds is corrected to
two seeds plus a hardcoded reference row. Retracting one's own claims in the
document that supersedes them is the minimum standard, not a virtue, and it is
recorded here so that any reader encountering the earlier artifacts knows their
status.

\paragraph{Environmental cost.} Not measured. [TODO: report GPU-hours for
training and for the transfer runs, and estimate the associated energy use.]

\paragraph{Dual use and population bias.} The cohort is critical-care patients
at one US academic medical centre in a historical period. The drug dictionary,
the class dictionary, and the danger judgements encode that context. Any system
built on this pipeline would inherit those biases, and its errors would not be
distributed uniformly across patient populations. Nothing in this work measured
subgroup performance, and with twelve positives nothing could.

\section{Summary}
\label{sec:eval-summary}

BioMedBERT fine-tuned on MedNLI reaches $83.90\%$ accuracy and $83.91\%$ macro
$F_1$, inside the published range. That result is a property of the backbone and
carries no ACCORD components; reporting it under the system's name would be the
most misleading available framing and is avoided.

On the $299$ pharmacist-annotated pairs, ACCORD predicts \emph{contradiction}
for $278$, achieving precision $0.047$ against a base rate of $0.040$, recall
$0.917$, and accuracy $0.115$. Four independent measurements establish that this
is not a threshold problem:

\begin{enumerate}
\item \textbf{Ranking.} AUROC $0.460$, stable across three treatments of the
      ambiguous class, average precision $0.039$ to $0.046$ against a $0.040$
      base rate, Mann-Whitney $p = 0.64$, bootstrap $95\%$ CI $[0.287, 0.633]$.
      The fifteen highest-scored pairs contain zero true contradictions.
\item \textbf{Operating points.} At thresholds $0.05$, $0.10$, and $0.15$ the
      model predicts positive for every pair at precision $0.0474$, exactly the
      base rate. Precision at recall $0.5$ is $5.4\%$.
\item \textbf{Severity.} Spearman $\rho = -0.0001$ ($p = 0.999$) between harm
      score and danger rating, from a head that never received a gradient. The
      severity research question was not tested.
\item \textbf{Comparison.} GPT-4o zero-shot is the only system with a positive
      MCC ($+0.109$). ACCORD is at chance ($-0.007$). The weak labeller it was
      distilled from is below chance ($-0.044$).
\end{enumerate}

Omission accounts for $7$ of the $12$ true contradictions, and it is the error
type an asymmetric cross-note encoder is structurally least able to represent,
because one side of the input is empty by construction. Four of the twelve are
rated $4$ or $5$ for danger, and no model score orders them.

The task is not impossible: a general-purpose model, zero-shot, on the same
evidence, performs above chance. What fails is this pipeline, and it fails at
every stage, for independent and individually diagnosable reasons.
